% =====================================================================
%  Chapitre 8 - Sprint 5 : facturation, alertes et journal metier
% =====================================================================
\chapter{Sprint 5 : facturation, alertes et journal metier}
\label{chap:sprint5}

\section*{Introduction}
\addcontentsline{toc}{section}{Introduction}

Les quatre sprints precedents ont produit de la matiere : des comptes,
des vehicules, des missions, des interventions. Ce sprint traite de trois
services qui n'en creent aucune : ils \essentiel{exploitent ce que les
autres ont produit}. La facturation transforme les missions en creances,
les alertes transforment les donnees en signaux, le journal metier
transforme les evenements en information de pilotage.

Cette position d'aval explique qu'ils aient ete traites ensemble, et
tardivement : chacun suppose que les modules amont existent deja.

% =====================================================================
\section{Planification du sprint}

\subsection{Objectifs}

\begin{enumerate}
  \item Etablir des factures a partir des missions realisees, sans
    double facturation possible.
  \item Suivre le cycle de vie d'une facture et identifier les impayes.
  \item Surveiller automatiquement les echeances, les stocks et les
    retards, et en alerter les responsables.
  \item Consigner les evenements metier dans un journal consultable.
  \item Diffuser les alertes critiques hors de l'application.
\end{enumerate}

\subsection{Carnet du sprint}

\begin{table}[H]
\centering
\caption{Carnet du sprint 5}
\label{tab:backlog-s5}
\small
\begin{tabularx}{\textwidth}{@{}cp{2.3cm}Xcc@{}}
\toprule
\textbf{N\textsuperscript{o}} & \textbf{Acteur} & \textbf{Besoin exprime} &
\textbf{Prio.} & \textbf{Pts} \\ \midrule
1 & Responsable & Etablir une facture depuis les missions facturables & Must & 8 \\
2 & Responsable & Suivre le cycle de vie d'une facture & Must & 5 \\
3 & Responsable & Identifier les factures impayees & Should & 3 \\
4 & Client & Consulter ses propres factures & Must & 3 \\
5 & Responsable & Consulter et filtrer les alertes metier & Must & 5 \\
6 & Responsable & Lancer les controles metier & Should & 5 \\
7 & Administrateur & Supprimer des alertes traitees & Should & 3 \\
8 & Responsable & Consulter le journal des evenements metier & Could & 5 \\
9 & Responsable & Recevoir une alerte critique hors application & Could & 5 \\
10 & Responsable & Editer un rapport quotidien de synthese & Could & 5 \\
\midrule
\multicolumn{4}{@{}r}{\textbf{Total}} & \textbf{47} \\
\bottomrule
\end{tabularx}
\end{table}

% =====================================================================
\section{Facturation}

\figReelle{finance-usecase.png}
  {Cas d'utilisation du service de facturation}
  {fig:uc-finance}{0.53}

\figReelle{finance-sequence.png}
  {Sequence : etablissement d'une facture}
  {fig:seq-finance}{0.61}

Le cas \exemple{etablir une facture} \essentiel{inclut} la verification
des missions facturables : la liste proposee au gestionnaire ne contient
que les missions terminees qui ne sont pas deja rattachees a une facture.
\essentiel{La double facturation est donc rendue impossible en amont},
plutot que detectee apres coup.

Le montant n'est pas saisi : il est calcule a partir du cout de la
mission, auquel s'ajoute la taxe sur la valeur ajoutee au taux en
vigueur. Les trois montants sont conserves separement, hors taxe, taxe et
toutes taxes comprises, ce qui evite de recalculer un historique si le
taux change.

Une facture parcourt cinq etats : brouillon, emise, payee, annulee et en
retard. Le passage en retard n'est pas une action de l'utilisateur mais
\essentiel{une consequence du temps} : il est deduit de la date
d'echeance, et c'est le service d'alertes qui le signale.

% =====================================================================
\section{Alertes metier}

\figReelle{alerts-usecase.png}
  {Cas d'utilisation du service d'alertes}
  {fig:uc-alerts}{0.69}

\figReelle{alerts-sequence.png}
  {Sequence : execution des controles metier}
  {fig:seq-alerts}{0.55}

Le service d'alertes ne detient aucune donnee propre : \essentiel{il lit
celles des autres domaines} et en tire des signaux. La
figure~\ref{fig:seq-alerts} montre le mecanisme : un balayage parcourt
les echeances documentaires, les assurances, les permis, les seuils de
stock et les missions en retard, puis produit une alerte pour chaque
situation rencontree.

\paragraph{Deux precautions de conception.}
La premiere est l'\essentiel{ecart des doublons} : le controle verifie
qu'une alerte identique et encore active n'existe pas deja avant d'en
creer une nouvelle. Sans cela, un vehicule dont la visite technique est
expiree produirait une alerte a chaque execution.

La seconde tient au \essentiel{vocabulaire}. Le type d'une alerte n'est
pas une chaine libre mais une valeur d'enumeration, propre au service.
Cette enumeration compte quarante et une valeurs et \essentiel{differe de
celle du projet partage} : ecrire une valeur du projet partage qui n'y
figure pas fait echouer la lecture. La contrainte a ete rencontree en
conditions reelles lors de la constitution du jeu de demonstration.

% =====================================================================
\section{Journal metier}

\figReelle{iboa-usecase.png}
  {Cas d'utilisation du service de journal metier}
  {fig:uc-iboa}{0.84}

\figReelle{iboa-communication.png}
  {Echanges du service de journal metier}
  {fig:com-iboa}{0.66}

Le service de journal metier occupe une position particuliere :
\essentiel{il ne consulte rien, il ecoute}. Ecrit en Python, il souscrit
a l'echange par sujet du courtier de messages et recoit les evenements
publies par les services metier. Pour chacun, il evalue une criticite,
regroupe les evenements lies et, au-dela d'un seuil, diffuse une alerte
par courriel ou par message instantane.

C'est le seul service de la plateforme dont \essentiel{le point d'entree
principal n'est pas une requete mais un message}. Son interface web ne
sert qu'a consulter ce qu'il a deja consigne.

\begin{table}[H]
\centering
\caption{Entites principales des trois services}
\label{tab:entites-s5}
\small
\begin{tabularx}{\textwidth}{@{}p{2.6cm}p{3.4cm}X@{}}
\toprule
\textbf{Service} & \textbf{Entite} & \textbf{Role et particularite} \\ \midrule
Facturation & Facture & Montants hors taxe, taxe et toutes taxes
  conserves separement ; rattachement optionnel a une mission \\
Facturation & Devis & Montants et validite ; distinct de la demande de
  transport du sprint~3, qu'il prolonge cote comptable \\
Alertes & Alerte & Type, priorite, statut, et references facultatives
  vers le vehicule, le chauffeur, l'equipement ou la mission concernes \\
Journal metier & Evenement & Charge utile brute recue du courtier,
  criticite evaluee, horodatage de reception \\
Journal metier & Alerte metier & Produite par regroupement d'evenements
  lies ; porte le canal de diffusion utilise \\
Journal metier & Preferences & Seuils et canaux choisis par le
  responsable pour la diffusion \\
\bottomrule
\end{tabularx}
\end{table}

% =====================================================================
\section{Realisation}

\subsection{La facturation}

La liste des factures, figure~\ref{fig:s5-factures}, est organisee autour
du \essentiel{statut de reglement} : brouillon, emise, payee, en retard,
annulee. C'est la seule lecture qui interesse reellement le service
administratif, dont la question quotidienne n'est pas \exemple{combien
avons-nous facture} mais \exemple{qu'est-ce qui n'est pas encore rentre}.
Les montants hors taxe, la taxe et le montant toutes taxes comprises sont
affiches separement, conformement au choix de modelisation decrit plus
haut : les recalculer a l'affichage aurait expose le document a des
divergences d'arrondi entre l'ecran et le document remis au client.

\capture{s5-liste-factures.jpg}
  {Factures et statut de reglement}
  {fig:s5-factures}

L'etablissement d'une facture, figure~\ref{fig:s5-creation}, ne part pas
d'une page blanche. L'ecran propose les \essentiel{missions facturables}
du client retenu, c'est-a-dire celles qui sont terminees et qui ne sont
rattachees a aucune facture. Le montant est preleve sur le devis accepte
au sprint~3. Cette continuite entre le devis et la facture est ce qui
supprime la double saisie que l'entreprise pratiquait auparavant entre son
tableur de suivi et son logiciel comptable, et avec elle les ecarts qu'une
recopie manuelle finit toujours par produire.

\capture{s5-creation-facture.jpg}
  {Etablissement d'une facture a partir des missions facturables}
  {fig:s5-creation}

\subsection{La surveillance}

Le centre d'alertes, figure~\ref{fig:s5-alertes}, rassemble en un seul
point ce qui reclame une decision : une assurance qui expire, une visite
technique depassee, un vehicule dont la revision est due, un equipement
au-dela de son seuil d'usure. Chaque alerte porte une priorite et un
statut, et peut etre rattachee au vehicule, au chauffeur, a l'equipement
ou a la mission concernes. Le filtre par priorite est ici l'element
determinant : sans lui, un centre d'alertes se remplit et cesse d'etre
consulte, ce qui le rend plus nuisible qu'utile.

\capture{s5-centre-alertes.jpg}
  {Centre d'alertes metier, filtrable par priorite et par statut}
  {fig:s5-alertes}

Le journal metier, figure~\ref{fig:s5-iboa}, releve d'une logique
differente. Il ne signale pas ce qu'il faut faire : il conserve ce qui
s'est passe. Chaque evenement publie par un service sur le courtier de
messages y est consigne avec sa criticite et son horodatage. L'ecran sert
a l'analyse apres coup --- comprendre pourquoi une mission a derape,
retrouver la chronologie d'un incident --- et non au pilotage courant.
C'est cette difference de finalite qui a justifie deux modules distincts
plutot qu'un seul.

\capture{s5-journal-iboa.jpg}
  {Journal des evenements metier et de leur criticite}
  {fig:s5-iboa}

% =====================================================================
\section{Tests de validation}

\begin{table}[H]
\centering
\caption{Tests de validation du sprint 5}
\label{tab:tests-s5}
\small
\begin{tabularx}{\textwidth}{@{}cXp{2.9cm}c@{}}
\toprule
\textbf{N\textsuperscript{o}} & \textbf{Cas teste} & \textbf{Attendu} &
\textbf{Etat} \\ \midrule
T1 & Mission deja facturee & Absente de la liste des facturables & Reussi \\
T2 & Facture etablie & Montants hors taxe, taxe et TTC coherents & Reussi \\
T3 & Echeance depassee sans reglement & Facture signalee en retard & Reussi \\
T4 & Controles metier lances deux fois & Aucune alerte en double & Reussi \\
T5 & Evenement publie sur le courtier & Consigne au journal avec sa criticite & Reussi \\
T6 & Alerte critique produite & Diffusion hors application declenchee & Reussi \\
\bottomrule
\end{tabularx}
\end{table}

% =====================================================================
\section*{Conclusion}
\addcontentsline{toc}{section}{Conclusion}

Ce sprint livre les trois services d'aval : facturation sans double
emploi possible, surveillance automatique des echeances et des seuils, et
journal metier alimente par le courtier de messages. Les quarante-sept
points du carnet ont ete realises.

Il met aussi en lumiere une \essentiel{limite du systeme d'evenements}.
La bibliotheque commune declare vingt-deux evenements d'integration ;
dans l'etat actuel, \essentiel{un seul service en publie et un seul en
consomme}. Le reste est declare sans emetteur ni destinataire. Ce n'est
pas un defaut de fonctionnement, mais une reserve de couplage qui n'a pas
ete exploitee faute de temps. Le chapitre de conclusion y revient au titre
des perspectives.

Le dernier sprint aborde ce qui reste : la collaboration entre les
utilisateurs et les deux services d'aide a la decision non encore
traites.
